\centerline{\bf \Huge Конструктивы I}


\begin{flushright}
        \scriptsize{Единственный счастливый предмет, который вам понадобится, уже находится внутри — УВЕРЕННОСТЬ\\ Мидорима Синтаро}
\end{flushright}

\problem{\textbf{1}} В клетках таблицы (4 строки, 7 столбцов) расставьте целые числа так, чтобы их сумма в каждой строке была равна 35, а в каждом столбце 20. Найдите хотя бы два решения.

\problem{\textbf{2}} В квадрате $3\times3$ расставьте числа 0, 1, 2, 3, 4, 5, 6, 7, 8 так, чтобы сумма чисел в каждой строке, столбце и диагонали была одинакова.

\problem{\textbf{3}} На доске написано число. За ход разрешается либо удвоить его, либо стереть у него последнюю цифру. Вначале на доске написано число 458. Как за несколько ходов получить число 14?

\problem{\textbf{4}} Разложите гири с весами 1, 2, ..., 555 на три кучки, равные по весу.

\problem{\textbf{5}} Можно ли на ребрах куба расставить числа от 1 до 12 так, чтобы все суммы чисел на гранях были одинаковы? (Конечно, можно - это же конструктивы)

\problem{\textbf{6}} Начальник должен посадить 8 учеников в 7 рядов по 3 ученика в каждом ряду. Как ему это сделать?

\problem{\textbf{7}} Тетрадный лист имеет размеры 33×41 клеток. Верно ли, что в его клетки можно записать все числа от 1 до 1353 так, чтобы для всех квадратиков размером 2×2 сумма записанных в них четырех чисел была одна и та же? (Сложная задача)